\chapter{$p$-adic analysis and formal group laws}
Throughout this chapter, $O$ is a complete discrete valuation ring with fraction field of
characteristic $0$, residue characteristic $p$, maximal ideal $\mathfrak m = (\pi)$ and
normalized valuation $v$. A natural number $e$ with $v(p) \le e$ bounds the absolute
ramification index; the unramified case is $e = 1$, where $\pi = p$. Nothing in this chapter
mentions the curve.

\begin{definition}[Evaluation of power series]
\label{def:eval}
\lean{ChabautyColeman.PSeries.eval, ChabautyColeman.PSeries.hasSum_eval,
  ChabautyColeman.PSeries.hasEval_of_mem, PadicInt.isAdic}
\leanok
For $h = \sum a_k z^k \in O[[z]]$ and $z_0 \in \mathfrak m$, set
$h(z_0) := \sum a_k z_0^k$; the series converges since $v(a_k z_0^k) \to \infty$.
\emph{Formalization:} a thin layer over Mathlib's power-series evaluation; the standing
setting is a complete Hausdorff uniform ring $O$, local with the $\mathfrak m$-adic
topology, so that elements of $\mathfrak m$ are topologically nilpotent. The
identification of the metric topology of $\ZZ_p$ with the $p\ZZ_p$-adic one makes
$O = \ZZ_p$ an instance.
\end{definition}

\begin{lemma}[Evaluation calculus]
\label{lem:eval-hom}
\lean{ChabautyColeman.PSeries.evalRingHom, ChabautyColeman.PSeries.comp,
  ChabautyColeman.PSeries.eval_comp, ChabautyColeman.PSeries.constantCoeff_comp,
  ChabautyColeman.PSeries.eval_deriv, ChabautyColeman.PSeries.summable_of_mem_pow}
\leanok
\uses{def:eval}
The map $h \mapsto h(z_0)$ is a ring homomorphism; evaluation commutes with composition
$h \circ k$ for $k(0) \in \mathfrak m$; and the formal derivative controls first
differences: $h(z_0 + w) = h(z_0) + h'(z_0)\,w + r\,w^2$ for some $r \in O$ and all
$z_0, w \in \mathfrak m$. (That a power series vanishing identically on $\mathfrak m$ is
$0$ when the residue field is infinite is the one-variable case of
Theorem~\ref{thm:interpolation}.)
\end{lemma}
\begin{proof}
\leanok
The ring-homomorphism and composition statements follow from continuity of evaluation
(the composition is defined by evaluating $h$ at the point $k$ of $O[[z]]$ for the
product topology). For the Taylor statement, peel the first two terms off the binomial
expansion of $(z_0 + w)^n$ termwise; the remainder series has terms in
$\mathfrak m^{n-2}$, hence converges.
\end{proof}

\begin{theorem}[Weierstrass preparation]
\label{thm:wprep}
\lean{PowerSeries.exists_isWeierstrassFactorization,
  PowerSeries.IsWeierstrassFactorization.elim,
  ChabautyColeman.PSeries.instIsAdicComplete}
\leanok
\uses{def:eval}
If $h \in O[[z]]$ has $\bar h \ne 0$ with $m := \ord(\bar h) < \infty$, then $h = u W$
with $u \in O[[z]]^\times$ and $W \in O[z]$ monic of degree $m$ with
$W \equiv z^m \pmod{\mathfrak m}$, and this factorization is unique.
\end{theorem}
\begin{proof}
\leanok
This is in Mathlib for local rings that are adically complete with respect to the
maximal ideal (distinguished polynomials are by definition monic with lower coefficients
in $\mathfrak m$, and their degree is $\ord(\bar h)$). What remains is that the standing
topological hypotheses imply $\mathfrak m$-adic completeness.
\end{proof}

\begin{corollary}[Zeros in the open unit disk]
\label{cor:wprep-count}
\lean{ChabautyColeman.PSeries.zeros_le_ord, ChabautyColeman.PSeries.finite_setOf_eval_eq_zero,
  ChabautyColeman.PSeries.setOf_eval_eq_zero_eq, ChabautyColeman.PSeries.isUnit_eval_of_isUnit}
\leanok
\uses{thm:wprep, lem:eval-hom}
For $O$ a domain, such an $h$ has at most $m$ distinct zeros in $\mathfrak m$; the zeros
are exactly the roots of $W$, all of which lie in $\mathfrak m$. Since the statement
applies to every $O$ satisfying the standing hypotheses, it applies in particular to
every complete extension $O'/O$. (A refinement counting zeros with multiplicity is
possible, but only the count of distinct zeros is used: rational points inject into
zeros.)
\end{corollary}
\begin{proof}
\leanok
Values of a power series on $\mathfrak m$ are congruent to the constant coefficient
modulo $\mathfrak m$ (write $h - h(0) = z\,g$ and evaluate), so units of $O[[z]]$ do
not vanish on $\mathfrak m$: $h$ and $W$ have the same zeros there. Roots of a
distinguished polynomial lie in $\mathfrak m$ ($z^m \in \mathfrak m$ at a root, and
$\mathfrak m$ is prime), and a monic polynomial of degree $m$ over a domain has at most
$m$ roots; $\deg W = m$ by comparison with the reduction.
\end{proof}

\begin{theorem}[Strassmann]
\label{thm:strassmann}
\lean{ChabautyColeman.PSeries.strassmann, ChabautyColeman.PSeries.strassmann_finite,
  ChabautyColeman.PSeries.evalT, ChabautyColeman.PSeries.recenter,
  ChabautyColeman.PSeries.evalT_add_eq, ChabautyColeman.PSeries.map_recenter}
\leanok
\uses{thm:wprep, cor:wprep-count, lem:eval-hom}
Let $O$ be a domain and $h = \sum a_k z^k \in O[[z]]$ with $a_k \to 0$ (so that $h$
converges on all of $O$), $a_N$ a unit, and $a_k \in \mathfrak m$ for all $k > N$. Then
$h$ has at most $N$ zeros in $O$. (The decay hypothesis $a_k \to 0$ is necessary for
convergence outside $\mathfrak m$; finiteness of the number of unit coefficients alone
does not suffice.)
\end{theorem}
\begin{proof}
\leanok
Instead of preparation at radius $1$, which would require the Tate algebra, decompose
the zero set by residue classes: the reduction $\bar h$ is a polynomial $q$ of degree
$N$; zeros of $h$ lie over roots of $q$; over a root $\bar c$ with lift $c$, the zeros
of $h$ are the zeros on $\mathfrak m$ of the recentred series $h(c + z)$ (defined
coefficientwise by convergent series; the identity between $h(c+w)$ and the recentring
evaluated at $w$ is an exchange of summation). The reduction of $h(c + z)$ is the
Taylor shift $q(\bar c + z)$, whose vanishing order is the multiplicity of $q$ at
$\bar c$, so Corollary~\ref{cor:wprep-count} bounds the zeros in the class by that
multiplicity, and $\sum_{\bar c} \mathrm{mult}_{\bar c}(q) \le \deg q = N$.
\end{proof}

\begin{lemma}[An index-valuation inequality]
\label{lem:padic-index}
\lean{Nat.add_mul_multiplicity_lt, Nat.mul_multiplicity_le_of_lt}
\leanok
Let $p > m + e + 1$ and $k + 1 > m + 1$. Then $(k+1) - e\,v_p(k+1) > m + 1$. (No primality
of $p$ is needed; in Lean the inequality is stated without natural subtraction, as
$m + 1 + e\,v_p(k+1) < k+1$.)
\end{lemma}
\begin{proof}
\leanok
Let $j := v_p(k+1)$. For $j = 0$ this is the hypothesis $k + 1 > m + 1$. For $j \ge 1$,
$k + 1 \ge p^j \ge (m+e+2)^j \ge (m+e+2)\,j \ge m + 2 + e\,j$, the middle step because
$c\,j \le c^j$ for $c \ge 2$ and the last because $(m+2)(j-1) \ge 0$.
\end{proof}

\begin{theorem}[Zeros of antiderivatives]
\label{thm:antideriv-zeros}
\lean{ChabautyColeman.PSeries.antideriv_zeros_le,
  ChabautyColeman.PSeries.antideriv_zeros_finite,
  ChabautyColeman.PSeries.pow_dvd_coeff_antideriv,
  IsLocalRing.IsUniformizer.exists_pow_eq,
  Nat.tendsto_sub_mul_multiplicity_atTop}
\leanok
\uses{thm:strassmann, lem:padic-index}
Let $\mathfrak m = (\pi)$ with $v(p) \le e$, let
$w = \sum_{k\ge0} b_k z^k \in O[[z]]$ with $b_m$ a unit, and let
$G(z) = c + \sum_k \tfrac{b_k}{k+1} z^{k+1}$ be an antiderivative, $c \in O$. If
$p > m + e + 1$, then $G$ has at most $m + 1$ zeros in $\mathfrak m$. ($m$ need not be
the exact reduction order of $w$: any index with unit coefficient works, and
$m = \ord(\bar w)$ gives the sharpest bound. For $v(c) < 0$ the statement is trivial, so
requiring $c \in O$ loses nothing.)
\end{theorem}
\begin{proof}
\leanok
\uses{lem:padic-index, thm:strassmann}
Work with the substituted series $H(u) := G(\pi u)$, which has coefficients in $O$ and is
characterized division-free by $(k+1)\,\mathrm{coeff}_{k+1}(H) = \pi^{k+1} b_k$; the claim
becomes a bound on its zeros in all of $O$. Integrality is the divisibility
$(k+1) \mid \pi^{k+1}$, which holds because $v(k+1) = v(p)\,v_p(k+1) \le e\,v_p(k+1)$ and
$(e+2)\,v_p(n) \le (e+2)^{v_p(n)} \le p^{v_p(n)} \le n$ for $p \ge e+2$. The coefficient of
$H$ at index $k+1$ has valuation $v(b_k) + (k+1) - v(k+1)$, hence at least
$v(b_k) + (k+1) - e\,v_p(k+1)$: it equals $m+1$ at index $m+1$ (using $m + 1 < p$, so
integers prime to $p$ are units in $O$), exceeds $m+1$ for every index $> m+1$ by
Lemma~\ref{lem:padic-index}, and the coefficients tend to $0$ (by the same estimate, since
$(e+2)\,v_p(k+1) \le k+1$). Dividing by $\pi^{s}$ with $s$ the minimal valuation ($\le m+1$)
gives a series with unit content and top unit index $\le m+1$; Strassmann's theorem
bounds its zeros in $O$ by $m+1$.
\end{proof}

\begin{definition}[Multivariate evaluation]
\label{def:mv-eval}
\lean{ChabautyColeman.MvPSeries.eval, ChabautyColeman.MvPSeries.hasEval_of_mem,
  ChabautyColeman.MvPSeries.evalRingHom, ChabautyColeman.MvPSeries.hasSum_eval,
  MvPowerSeries.continuous_subst', ChabautyColeman.MvPSeries.eval_subst,
  ChabautyColeman.MvPSeries.hasEval_eval}
\leanok
\uses{def:eval}
For $h \in O[[z_1,\dots,z_n]]$ and $\mathbf z_0 \in \mathfrak m^n$, the value
$h(\mathbf z_0)$ is defined by the same estimates. Evaluation is a ring homomorphism and
is compatible with formal substitution: substitution is continuous for the product
topologies (each coefficient of a substitution is a fixed finite linear combination of
the input coefficients), so evaluating a substitution equals evaluating at the evaluated
family, which is itself a legitimate evaluation point.
\end{definition}

\begin{definition}[Evaluation of series with decaying coefficients]
\label{def:mv-evalT}
\lean{ChabautyColeman.MvPSeries.evalT, ChabautyColeman.MvPSeries.hasSum_evalT,
  ChabautyColeman.MvPSeries.evalT_eq_eval, ChabautyColeman.MvPSeries.evalT_add,
  ChabautyColeman.MvPSeries.tendsto_coeff_zero_of_mem_pow,
  Finsupp.tendsto_degree_cofinite}
\leanok
\uses{def:mv-eval}
For $h \in O[[z_1,\dots,z_n]]$ whose coefficients tend to $0$ (along the cofinite filter
on multi-indices) and \emph{any} family $\mathbf z \in O^n$ --- not restricted to
$\mathfrak m^n$ --- the sum $h(\mathbf z) = \sum_d (\mathrm{coeff}_d h)\,\mathbf z^d$
converges, since each term lies in a power of $\mathfrak m$ determined by
$\mathrm{coeff}_d h$, the point family being bounded. This total evaluation is additive
in $h$ and agrees with the evaluation of Definition~\ref{def:mv-eval} wherever the
latter converges. A convenient sufficient condition for the decay hypothesis:
$\mathrm{coeff}_d h \in \mathfrak m^{e(d)}$ with $e(d) \to \infty$ cofinitely, which
holds once $e$ grows with the total degree, since only finitely many multi-indices have
bounded degree. This is the $n$-variable analogue of the coefficient-decay evaluation
behind Strassmann's theorem; it is what evaluates the logarithm on~$\mathfrak m^n$ after
scaling by~$\pi$.
\end{definition}

\begin{theorem}[Interpolation]
\label{thm:interpolation}
\lean{ChabautyColeman.MvPSeries.eq_zero_of_forall_eval_zero,
  ChabautyColeman.PSeries.eq_zero_of_forall_eval_zero}
\leanok
\uses{def:mv-eval, cor:wprep-count}
Let $O$ be a domain of the standing setting with $\mathfrak m = (\pi)$ principal,
$\pi \ne 0$, and $O$ infinite. If $h \in O[[z_1,\dots,z_n]]$ vanishes at every point of
$\mathfrak m^n$, then $h = 0$. No condition on the residue field is needed; in
particular the statement applies directly over $\ZZ_p$.
\end{theorem}
\begin{proof}
\leanok
One variable: if the reduction of $h$ were nonzero, the zero set of $h$ in $\mathfrak m$
would be finite (Corollary~\ref{cor:wprep-count}), but it is all of the infinite set
$\mathfrak m$; hence $\pi$ divides all coefficients, and iterating, all coefficients lie
in $\bigcap_s \mathfrak m^s = 0$. Several variables, by induction on $n$: fixing all but
one variable, the coefficient slices collapse into a one-variable series computing the
evaluation (an exchange of summation), and the one-variable case applies
coefficientwise.
\end{proof}

\begin{definition}[Formal group law]
\label{def:fgl}
\lean{ChabautyColeman.FormalGroupLaw, ChabautyColeman.FormalGroupLaw.Points}
\leanok
\uses{def:mv-eval}
An $n$-dimensional commutative formal group law over $O$ is a tuple
$F \in O[[X;Y]]^n$ with $F(X,0)=X$, $F(0,Y)=Y$, $F(X,Y)=F(Y,X)$ and
$F(F(X,Y),Z)=F(X,F(Y,Z))$, as formal substitution identities. It makes $\mathfrak m^n$ a
commutative \emph{monoid} via Definition~\ref{def:mv-eval}: each monoid axiom is the
corresponding formal axiom evaluated, using that evaluation of a substitution is
evaluation at the evaluated family, and that values of series without constant term stay
in $\mathfrak m$. (Inverses, which would make this monoid a group, are never needed: the
group laws below parametrize residue polydisks of a group and inherit inverses from it.)
\end{definition}

\begin{remark}[Formalization]
The monoid of points (\texttt{ChabautyColeman.FormalGroupLaw.Points}) is formalized in
\texttt{HyperellipticCurves.Archive.FormalGroupPoints}; the main theorems do not depend on it,
since the main development works with the evaluated identities directly, in the $\pi$-scaled
form of Theorem~\ref{thm:log-iso}.
\end{remark}

\begin{lemma}[Formal partial derivatives]
\label{lem:pderiv}
\lean{MvPowerSeries.pderiv, MvPowerSeries.pderiv_subst, MvPowerSeries.continuous_subst'}
\leanok
Formal partial derivatives of multivariate power series satisfy the Leibniz rule and the
chain rule for formal substitution.
\end{lemma}
\begin{proof}
\leanok
The Leibniz rule holds by construction. For the chain rule, both sides are continuous for the
product topologies (with discrete coefficients) and agree on the dense subring of polynomials,
where the identity follows by induction from the polynomial chain rule.
\end{proof}

\begin{remark}[Formalization]
The partial derivative itself is Mathlib's \texttt{MvPowerSeries.pderiv}, which is a derivation,
so the Leibniz rule is \texttt{Derivation.leibniz} and needs no separate statement; the project
adds the chain rule for substitution and the continuity of substitution over an arbitrary
topological coefficient ring.
\end{remark}

\begin{definition}[Differential matrix]
\label{def:fgl-diff}
\lean{ChabautyColeman.FormalGroupLaw.diffMatrix,
  ChabautyColeman.FormalGroupLaw.diffMatrix_map_constantCoeff,
  ChabautyColeman.FormalGroupLaw.isUnit_diffMatrix}
\leanok
\uses{def:fgl, lem:pderiv}
$M(X) := (\partial F/\partial Y)(X,0) \in \mathrm{GL}_n(O[[X]])$: applying the chain
rule to $F(0,Y) = Y$ gives $M(0) = I$, so $\det M$ has unit constant term and $M$ is
invertible over the local ring $O[[X]]$.
\end{definition}

\begin{definition}[Base change of a formal group law]
\label{def:fgl-basechange}
\lean{ChabautyColeman.FormalGroupLaw.map,
  ChabautyColeman.FormalGroupLaw.diffMatrix_map,
  ChabautyColeman.FormalGroupLaw.diffMatrixInv_map}
\leanok
\uses{def:fgl, def:fgl-diff}
For a ring homomorphism $f \colon O \to O'$, applying $f$ to the coefficients of each
$F_i$ gives a formal group law $f_*\Phi$ over $O'$. The differential matrix and its
inverse commute with base change: $M(f_*\Phi) = f_*(M(\Phi))$ and
$M(f_*\Phi)^{-1} = f_*(M(\Phi)^{-1})$. In particular, since $M(\Phi)^{-1}$ has entries
in $O[[X]]$ already over $O$, the entries of the invariant differential $M^{-1}$ of the
base change to $\mathrm{Frac}\,O$ are integral; only the logarithm below acquires
denominators.
\end{definition}

\begin{definition}[Invariant differential and logarithm]
\label{def:fgl-log}
\lean{ChabautyColeman.FormalGroupLaw.log, ChabautyColeman.eulerAntideriv,
  ChabautyColeman.FormalGroupLaw.coeff_single_log,
  ChabautyColeman.FormalGroupLaw.diffMatrixInv_map_constantCoeff}
\leanok
\uses{def:fgl, def:fgl-diff}
The invariant differential is $\delta := M(X)^{-1}dX$. Over a field of coefficients, the
logarithm $\Lam$ is \emph{defined} as the radial (Euler) antiderivative of
$\sum_j (M^{-1})_{ij} X_j$: divide the degree-$k$ homogeneous part by $k$. Then
$\Lam \in X + (\deg \ge 2)$ and $\Lam(0) = 0$, the normalization holding because
$M^{-1}(0) = I$. That $d\Lam = \delta$ --- for which $\Lam$ is then the unique solution
with $\Lam(0) = 0$, and which requires the closedness of $\delta$ --- is
Lemma~\ref{lem:fgl-log-diff}.
\end{definition}

\begin{lemma}[Matrix invariance identity]
\label{lem:fgl-invariance-matrix}
\lean{ChabautyColeman.FormalGroupLaw.diffMatrix_subst_F}
\leanok
\uses{def:fgl-diff, lem:pderiv}
$M(F(X,Y)) = (D_Y F)(X,Y)\cdot M(Y)$ as matrices over $O[[X,Y]]$, where
$(D_Y F)_{ij} = \partial F_i/\partial Y_j$ and $M(Y)$ is the differential matrix in the
variables of the second block.
\end{lemma}
\begin{proof}
\leanok
Differentiate the associativity law $F(F(X,Y),C) = F(X,F(Y,C))$ with respect to $C_k$
and set $C = 0$. By the chain rule (Lemma~\ref{lem:pderiv}), the left side becomes
$M(F(X,Y))_{ik}$ and the right side $\sum_j \partial_{Y_j}F_i(X,Y)\cdot M(Y)_{jk}$,
using $F(Y,0) = Y$ to collapse the inner substitutions.
\end{proof}

\begin{lemma}[Closedness of the invariant differential]
\label{lem:fgl-invariance-closed}
\lean{ChabautyColeman.FormalGroupLaw.pderiv_diffMatrixInv_symm}
\leanok
\uses{def:fgl-diff, lem:fgl-invariance-matrix, lem:pderiv}
The invariant differential $\delta_i = \sum_a (M^{-1})_{ia}\,dX_a$ is closed:
$\partial_l (M^{-1})_{ik} = \partial_k (M^{-1})_{il}$. Commutativity of $F$ is
essential.
\end{lemma}
\begin{proof}
\leanok
Inverting the matrix identity of Lemma~\ref{lem:fgl-invariance-matrix} gives
$N(F(X,Y))\,(D_Y F)(X,Y) = N(Y)$ for $N = M^{-1}$. Differentiate with respect to $Y_l$
and antisymmetrize in $(k,l)$: the second-derivative terms cancel since partial
derivatives commute, leaving, with $A_{iab} = \partial_b N_{ia} - \partial_a N_{ib}$,
\[ A_{ikl}(Y) = \sum_{a,b} A_{iab}(F(X,Y))\,\partial_{Y_k}F_a\,\partial_{Y_l}F_b. \]
Setting $Y = 0$ (so $F(X,0) = X$ and $\partial_{Y_k}F_a(X,0) = M_{ak}$) shows
$M^{\mathsf T} A_i M = A_i(0)$ is the \emph{constant} matrix of linear coefficients.
From $M^{-1}M = 1$ and $M(0) = M^{-1}(0) = 1$, the linear coefficients of $M^{-1}$ and
$M$ are opposite, and those of $M$ are the mixed second derivatives
$\partial_{X_l}\partial_{Y_k} F_i(0)$, symmetric in $(k,l)$ by commutativity of $F$
(applied as a renaming of variables). Hence $A_i(0) = 0$, and multiplying by the
inverses of $M^{\mathsf T}$ and $M$ gives $A_i = 0$.
\end{proof}

\begin{lemma}[Differential of the logarithm]
\label{lem:fgl-log-diff}
\lean{ChabautyColeman.FormalGroupLaw.pderiv_log}
\leanok
\uses{def:fgl-log, lem:fgl-invariance-closed, lem:pderiv}
$d\Lam = \delta$, i.e.\ $\partial_k \Lam_i = (M^{-1})_{ik}$, in characteristic zero.
\end{lemma}
\begin{proof}
\leanok
For a closed form, the radial antiderivative is a primitive: with the Euler operator
$E = \sum_j X_j\partial_j$, which acts on the degree-$n$ part as multiplication by $n$,
one has $E\Lam_i = \sum_j (M^{-1})_{ij}X_j$ by construction, and applying $\partial_k$
gives $(E+1)(\partial_k\Lam_i) = \partial_k\sum_j (M^{-1})_{ij}X_j
= (E+1)\bigl((M^{-1})_{ik}\bigr)$, the last equality by closedness
(Lemma~\ref{lem:fgl-invariance-closed}). Since $E + 1$ multiplies the degree-$n$ part by
$n + 1 \ne 0$, it is injective on coefficients in characteristic zero.
\end{proof}

\begin{lemma}[Invariance and additivity]
\label{lem:fgl-invariance}
\lean{ChabautyColeman.FormalGroupLaw.diffMatrixInv_subst_F,
  ChabautyColeman.FormalGroupLaw.log_subst_F}
\leanok
\uses{def:fgl-log, lem:fgl-invariance-matrix, lem:fgl-log-diff}
$\delta(F(X,Y))\cdot d_Y F(X,Y) = \delta(Y)$ as formal identities; consequently
$\Lam(F(X,Y)) = \Lam(X) + \Lam(Y)$.
\end{lemma}
\begin{proof}
\leanok
Inverting the matrix identity of Lemma~\ref{lem:fgl-invariance-matrix} gives
$M(F(X,Y))^{-1}\,(D_Y F)(X,Y) = M(Y)^{-1}$, which is the first identity. For
additivity, consider $G := \Lam(F(X,Y)) - \Lam(X) - \Lam(Y)$. Its $Y$-partials vanish:
by the chain rule and $d\Lam = \delta$ (Lemma~\ref{lem:fgl-log-diff}), both
$\partial_{Y_k}\Lam(F(X,Y))$ and $\partial_{Y_k}\Lam(Y)$ equal the $Y$-substituted
entry $(M^{-1})_{ik}(Y)$ by the first identity, while $\Lam(X)$ has no $Y$-dependence.
The $X$-partials vanish by the same computation reflected through commutativity of $F$
(under the block swap, the first identity turns into its $X$-block counterpart
$M(F(X,Y))^{-1}\,(D_X F)(X,Y) = M(X)^{-1}$). Finally $G$ has zero constant term, and in
characteristic zero a series with vanishing constant term and vanishing first partials
is zero.
\end{proof}

\begin{lemma}[Integral coefficients of the logarithm]
\label{lem:log-coeff-integral}
\lean{ChabautyColeman.FormalGroupLaw.coeff_log_map}
\leanok
\uses{def:fgl-log, def:fgl-basechange}
For $\Phi$ over $O$ and any homomorphism $f \colon O \to K$ into a field, the
coefficient of the logarithm of the base change $f_*\Phi$ at multi-index $d$ is
$(\deg d)^{-1}\, f(c_d)$, where $c_d \in O$ is the corresponding coefficient of the
integral row $\sum_j (M^{-1})_{ij} X_j$ of the invariant differential. Thus the
numerator is integral and the only denominator is the total degree $\deg d$.
\end{lemma}
\begin{proof}
\leanok
The differential matrix and its inverse commute with base change
(Definition~\ref{def:fgl-basechange}), so the row $\sum_j (M^{-1})_{ij} X_j$ of
$f_*\Phi$ is $f$ applied coefficientwise to that of $\Phi$; the Euler antiderivative
then divides the degree-$d$ part by $\deg d$.
\end{proof}

\begin{lemma}[$p$-power denominators of the logarithm]
\label{lem:log-coeff-padic}
\lean{ChabautyColeman.FormalGroupLaw.coeff_log_map_padic}
\leanok
\uses{lem:log-coeff-integral}
Over a local ring $O$ of residue characteristic $p$ (so that every integer prime to $p$
is a unit), the degree-$d$ coefficient of the logarithm of $f_*\Phi$ becomes integral
after multiplication by $f(p)^{v_p(\deg d)}$: the denominator is a pure power of $p$,
with exponent $v_p(\deg d)$.
\end{lemma}
\begin{proof}
\leanok
By Lemma~\ref{lem:log-coeff-integral} the denominator is $\deg d$; writing
$\deg d = p^{v_p(\deg d)}\cdot m$ with $p \nmid m$, the prime-to-$p$ factor $m$ is a
unit in $O$, so only the $p$-power part survives as a denominator.
\end{proof}

\begin{lemma}[Integrality of the scaled logarithm]
\label{lem:log-scale-integral}
\lean{ChabautyColeman.FormalGroupLaw.coeff_log_map_scale}
\leanok
\uses{lem:log-coeff-padic}
Assume $p > e + 1$. Scaling the variables by the uniformizer $\pi$ makes the logarithm
integral: the degree-$d$ coefficient of $\Lam(\pi U)$ is $f(c)$ for some
$c \in \mathfrak m^{\deg d - e\,v_p(\deg d)}$. So $\Lam(\pi U) \in O[[U]]$, and since
$\deg d - e\,v_p(\deg d) \to \infty$ the coefficients tend to $0$, which drives the
convergence of $\Lam$ on $\mathfrak m^n$.
\end{lemma}
\begin{proof}
\leanok
Multiply the $p$-power denominator bound of Lemma~\ref{lem:log-coeff-padic} by
$\pi^{\deg d}$ and split the denominator off the numerator: writing $v := v_p(\deg d)$, we
have $e\,v \le \deg d$ because $(e+2)\,v \le (e+2)^v \le p^v \le \deg d$, so
$\pi^{\deg d} = p^{v}\,c'$ with $c' \in \mathfrak m^{\deg d - e v}$ --- the ramification
bound is spent exactly here.
\end{proof}

\begin{lemma}[Convergence of the logarithm on $\mathfrak m^n$]
\label{lem:log-converge}
\lean{ChabautyColeman.FormalGroupLaw.exists_integral_scaled_log}
\leanok
\uses{lem:log-scale-integral, def:mv-evalT}
Over a local ring $O$ of residue characteristic $p$ with $p > e + 1$: after scaling the
variables by the uniformizer $\pi$, the logarithm of $f_*\Phi$ descends to an integral
power series $H \in O[[U]]$ --- its base change is the $\pi$-scaled logarithm,
$f(\mathrm{coeff}_d H) = f(\pi)^{\deg d}\,\mathrm{coeff}_d \Lam$ --- whose coefficients
tend to $0$. Hence $\Lam$ converges on $\mathfrak m^n$, its value at
$\mathbf x = \pi\mathbf y$ being $f$ of the total evaluation of $H$ at $\mathbf y$, with
$v(\Lam(\mathbf x)) \ge 1$.
\end{lemma}
\begin{proof}
\leanok
By Lemma~\ref{lem:log-scale-integral} the degree-$d$ coefficient of $\Lam(\pi U)$ is
$f(c_d)$ for some $c_d \in \mathfrak m^{\deg d - e\,v_p(\deg d)}$; take $H$ with these
integral coefficients. Since $\deg d - e\,v_p(\deg d) \to \infty$ (the $p$-adic valuation
is sublinear and only finitely many multi-indices have bounded degree), the coefficients
tend to~$0$ and $H$ converges by Definition~\ref{def:mv-evalT}.
\end{proof}

\begin{lemma}[The formal exponential]
\label{lem:fgl-exp}
\lean{ChabautyColeman.FormalGroupLaw.exp, ChabautyColeman.FormalGroupLaw.log_comp_exp,
  ChabautyColeman.FormalGroupLaw.exp_comp_log}
\leanok
\uses{def:fgl-log, def:fgl}
The formal exponential $\mathrm{Exp}$ is the compositional inverse of $\Lam$:
$\Lam \circ \mathrm{Exp} = \mathrm{Exp} \circ \Lam = X$ as formal power series. Since
$\Lam \in X + (\deg \ge 2)$ has identity linear part, this inverse exists and is unique.
\end{lemma}
\begin{proof}
\leanok
The compositional inverse of a formal self-map with zero constant term and identity
linear part exists (by the fixed-point iteration $E^{k+1} = X - (\Lam - X)(E^k)$, which
stabilizes one coefficient degree per step because $\Lam - X$ has order $\ge 2$), is
two-sided, and is unique. Apply this to $\Lam$, whose linear part is the identity
because $M^{-1}(0) = I$.
\end{proof}

\begin{lemma}[Convergence of the exponential for odd $p$]
\label{lem:exp-converge}
\lean{ChabautyColeman.FormalGroupLaw.exp_converges, ChabautyColeman.FormalGroupLaw.exp_bounded,
  ChabautyColeman.Dexp}
\leanok
\uses{lem:fgl-exp, lem:log-converge}
The inverse series $\mathrm{Exp}$ of $\Lam$ converges on~$(pO)^n$, by the classical
estimate $v_p(k!) = (k - s(k))/(p-1)$, where $s(k)$ is the sum of the base-$p$ digits
of~$k$, and $p \ge 3$: the degree-$d$ coefficient of $\mathrm{Exp}$ becomes integral after
multiplication by $p^{\lfloor (\deg d - 1)/(p-1) \rfloor}$, and $\mathrm{Exp}(pU)$
descends to an integral power series with coefficients tending to~$0$, as for the
logarithm in Lemma~\ref{lem:log-converge}.
\end{lemma}
\begin{proof}
\leanok
The coefficient bound is proved by strong induction on $\deg d$ through the recursion
$\mathrm{Exp} = X - (\Lam - X)\circ\mathrm{Exp}$: the product slack from
$\prod_i \mathrm{Exp}_i^{m_i}$ cancels the logarithm's $1/\deg m$ denominator, and
$(p-1)\,v_p(m) \le m - 1$ closes the estimate. Convergence needs $p \ge 3$ via
$\deg d - (\deg d - 1)/(p-1) \to \infty$.
\end{proof}

\begin{remark}[Formalization]
The exponential (Lemmas~\ref{lem:fgl-exp} and~\ref{lem:exp-converge}) serves only the
injectivity of the logarithm. All three are formalized, in
\texttt{HyperellipticCurves.Archive.FormalGroupExp},
\texttt{HyperellipticCurves.Archive.ExpConverge} and
\texttt{HyperellipticCurves.Archive.LogInjective}; the main theorems do not depend on them.
\end{remark}

\begin{theorem}[The logarithm as a homomorphism]
\label{thm:log-iso}
\lean{ChabautyColeman.FormalGroupLaw.log_addHom, ChabautyColeman.FormalGroupLaw.scaledF,
  ChabautyColeman.FormalGroupLaw.scaledLog,
  ChabautyColeman.FormalGroupLaw.subst_scaledF_descent,
  ChabautyColeman.MvPSeries.evalT_subst}
\leanok
\uses{lem:fgl-invariance, lem:log-converge, def:mv-eval}
For $p > e + 1$, $\Lam \colon \mathcal G(F)(\mathfrak m) \to (\mathrm{Frac}\,O)^n$ is a
group homomorphism. ($\Lam$ is moreover injective --- indeed, in the unramified case, a
bijection onto $(pO)^n$ via $\mathrm{Exp}$ --- but this is never used.)
\end{theorem}
\begin{proof}
\leanok
Descend the formal additivity $\Lam(F(\mathbf X,\mathbf Y)) = \Lam(\mathbf X) +
\Lam(\mathbf Y)$ (Lemma~\ref{lem:fgl-invariance}) to the $\pi$-scaled integral forms and
evaluate. \emph{Formalization.} The value of $\Lam$ is realized as the total evaluation
(Definition~\ref{def:mv-evalT}) of the $\pi$-scaled logarithm $H$ (\texttt{scaledLog}), and
the $\pi$-scaled group sum $F(\pi\mathbf X, \pi\mathbf Y)/\pi$ (\texttt{scaledF}, integral
with decaying coefficients for any scaling factor in $\mathfrak m$) is turned into addition
(\texttt{log\_addHom}): the additivity identity descends along the $\pi$-rescaling
substitution --- injectivity of the coefficient map and associativity of substitution --- to
$H\circ F_\pi = H(\mathbf X) + H(\mathbf Y)$ for the scaled group law $F_\pi$
(\texttt{subst\_scaledF\_descent}), and total evaluation commutes with substitution
(\texttt{evalT\_subst}).
\end{proof}

\begin{theorem}[Hensel's lemma, smooth-point form]
\label{thm:hensel}
\lean{ChabautyColeman.exists_smooth_root, ChabautyColeman.hensel_unique,
  ChabautyColeman.exists_weierstrass_root, IsAdicComplete.henselianLocalRing}
\leanok
Let $g \in O[X]$ (not necessarily monic) with $g(a_0) \in \mathfrak m$ and
$g'(a_0) \in O^\times$. Then $g$ has a unique root $x \equiv a_0 \pmod{\mathfrak m}$. In
particular, for $f$ with $f(\omega) = 0$, $f'(\omega) \in O^\times$ and any
$y \in \mathfrak m$, there is a unique $x \equiv \omega$ with $f(x) = y^2$ --- the
implicit-function step of the disk charts. This is the full Hensel content the proof
consumes: every chart root is a single one-variable Hensel root, over $O$ or over the
$z$-adically complete base $O[[z]]$.
\end{theorem}
\begin{proof}
\leanok
Newton iteration $x \mapsto x - c\, g(x)$ with $c\,g'(a_0) = 1$, contracting in the
$\mathfrak m$-adic metric: the residual after $k$ steps lies in $\mathfrak m^{k+1}$, the
increments telescope to a limit by nonarchimedean summability, and the limit is a root
by Hausdorffness. Uniqueness: the difference of two roots is killed by a unit times
itself modulo $\mathfrak m$. \emph{Formalization.} For monic polynomials this is
Mathlib's Henselian-ring property: a local ring that is adically complete at its maximal
ideal is Henselian (\texttt{IsAdicComplete.henselianLocalRing}); a multivariate smooth-point
form is formalized in \texttt{HyperellipticCurves.Archive.HenselSmooth}, and the main theorems
do not depend on it.
\end{proof}
